\documentclass[12pt,a4paper]{article}
\usepackage[english]{babel}
\usepackage[margin=1in]{geometry}
\usepackage{setspace}
\usepackage[authoryear]{natbib}
\usepackage[colorlinks=true,linkcolor=blue,citecolor=blue,urlcolor=blue,backref=page]{hyperref}
\usepackage{amsmath}
\usepackage{amssymb}
\usepackage{booktabs}

\onehalfspacing
\setcitestyle{authoryear,open={(},close={)}}

\title{Institutional Choice in Postwar Japan:\\
Welfare Loans vs. Public Pawnshops, 1961}

\author{[Author Name]}
\date{August 2024}

\begin{document}

\maketitle

\begin{abstract}
How did the formation of the welfare state transform the provision of credit to low-income households? This paper examines institutional allocation when publicly funded welfare lending entered the credit market for low-income households in postwar Japan. Using the 1961 Kanagawa Survey of the Borderline Poor covering 6,152 households, I compare welfare-loan users and public-pawnshop users. The central finding is that welfare loans were not distributed uniformly to the poorest households, but concentrated among those with specific, addressable financing needs and economic foundations suggesting repayment capacity. The welfare state thus created a new boundary of creditworthiness within the low-income population. Three main results emerge: (1) business failure increased welfare-loan probability by 8.6 percentage points, prolonged illness by 5.1 percentage points, indicating purpose-specific targeting; (2) public-assistance history decreased welfare-loan probability by 15.6 percentage points while household assets increased it by 3.8 percentage points per item, showing allocation by household economic position rather than shock type alone; and (3) welfare lending did not displace older borrowing—households layered multiple credit sources, with previous pawnshop users 39.1 percentage points more likely to report pawning when expenses were insufficient. The results reveal that state entry into credit markets changed not only credit availability but the criteria by which low-income households were sorted across institutions.
\end{abstract}

\newpage

\section{Introduction}

How did the formation of the welfare state transform the provision of credit to low-income households? Before the development of modern social protection, working-class households in Britain and Japan managed income shocks through combined use of precautionary savings, adjustments in family labor supply, informal borrowing, purchases on account, and pawn lending \citep{ScottWalker2012, Ogasawara2024}. Pawn credit served as a particularly important risk-coping mechanism for financially excluded populations. During the influenza pandemic of 1918--20, Japanese pawnshop lending increased as established users intensified borrowing in response to health shocks \citep{Inoue2021}. The financial lives of low-income households cannot therefore be understood through conventional bank credit alone. Instead, they depended on a portfolio of institutions, each serving distinct purposes and accessible under different conditions.

Much less is known about how this financial ecology changed once the state began lending directly to low-income households. Scholars have documented how publicly supported pawn institutions in Europe—the French \emph{Monts-de-Piété} in Bordeaux and failed experiments in Ireland—attempted to supply affordable credit to borrowers excluded from conventional banking \citep{McLaughlin2013, PastureaBlancheton2014}. Comparative work on the United States and France has traced how private consumer credit came to be justified as an instrument of social welfare and inclusion \citep{Trumbull2012}. Yet little is known about the allocation of households across multiple credit institutions when the welfare state itself became a lender to the poor. Japan provides an informative case. After 1955, the Household Rehabilitation Fund offered public loans for business, medical care, housing, and education to low-income households unable to access ordinary market credit. These loans operated alongside public pawnshops, informal borrowing, and consumption compression. The question remains: did welfare lending simply displace pawnshops and other credit sources, or did it introduce new criteria for distinguishing which low-income households could access which institutions?

This paper explores how low-income households were allocated across financial institutions when the welfare state entered the credit market in postwar Japan. Specifically, it examines whether welfare lending displaced public pawnshops or instead added a new institution serving a distinct group of households. The central finding is that welfare loans were not distributed uniformly to the poorest households. Rather, they were concentrated among households with identifiable financing needs and with an economic and material basis that could make rehabilitation and repayment appear feasible to welfare officials. The formation of the welfare state therefore did not simply incorporate low-income households into formal finance. It created a new boundary of creditworthiness within the low-income population by introducing a new lending technology based on assessments of future household rehabilitation.

The analysis draws on the 1961 \emph{Kanagawa Survey of the Borderline Poor}, a restored dataset covering 6,152 low-income households in six cities. The survey records the causes of low income—unemployment, business failure, declining earnings, and illness—alongside histories of public assistance, four types of welfare loans, and public-pawnshop use. It also contains information on income, employment, household assets, and coping methods when living expenses were insufficient. These features enable comparison of welfare-loan users and public-pawnshop users within the same low-income population. The survey records institutional histories but not loan dates, loan amounts, applications, rejections, or household circumstances before borrowing, so the analysis describes patterns of institutional sorting rather than causal effects of credit receipt.

The results yield three main findings. First, welfare lending was associated with specific, addressable problems. Even after controlling for income, location, household composition, and employment, business failure is associated with an 8.6-percentage-point higher probability of welfare-loan use, while prolonged illness is associated with a 5.2-percentage-point increase. The individual programmes also correspond closely to their stated purposes: Household Rehabilitation Fund use is associated with self-employment and business capital demand; educational lending with schooling needs; and medical lending with treatment needs. Welfare lending therefore operated less as unrestricted support for ordinary living expenses than as purpose-specific credit tied to particular plans for household rehabilitation.

Second, the clearest distinction between welfare lending and public pawnshops lay not in shock type but in household economic position. Among households using either welfare loans or pawnshops but not both, a public-assistance history is associated with a 15.6-percentage-point lower probability of belonging to the welfare-loan group. Households recorded as having low livelihood capacity are 14.5 percentage points less likely to use welfare loans, and households whose principal breadwinner was in poor health are 11.4 percentage points less likely. Conversely, each additional household asset is associated with an approximately 3.8-percentage-point higher probability of welfare-loan rather than pawnshop use. Since business failure is associated with both institutions, their division of labor cannot be reduced to shock-type separation alone. Instead, the pattern is consistent with welfare lending depending on household prospective capacity for rehabilitation, whereas public pawnshops supplied credit against currently available physical assets.

Third, welfare lending did not fully replace older forms of borrowing. Welfare-loan users were more likely to borrow from friends and neighbors; households with prior public-pawnshop use remained significantly more likely to report pawning goods when living expenses were insufficient—approximately 38 percentage points more likely. These associations are not causal, but they indicate that postwar financial inclusion did not involve transition to a single institution. Instead, households continued to layer multiple credit sources.

This paper contributes to the economic history of household risk coping and alternative finance. Historians have established that working-class households combined formal and informal credit, savings, and labor-supply adjustments before mature welfare states emerged \citep{ScottWalker2012, Ogasawara2024, Inoue2021}. Public pawn institutions in Europe attempted to extend affordable credit to the financially excluded, although success depended on regulation, local demand, and institutional design \citep{McLaughlin2013, PastureaBlancheton2014}. This literature primarily examines how individual institutions functioned before or outside mature welfare states. By comparing users of collateral-based finance and welfare lending within the same low-income population, this paper demonstrates that the state's entry into credit markets altered not merely credit availability but the institutional criteria by which borrowers were sorted.

Second, the paper adds a historical perspective to international scholarship on credit and welfare. Existing work interprets credit access both as an instrument of social inclusion and as a substitute for conventional welfare provision, with particular attention to how private consumer credit became justified as social welfare in the United States and France, and how credit-based policies figure in contemporary welfare regimes \citep{Trumbull2012, Mertens2017}. This literature emphasizes mature welfare states and private credit expansion in the later twentieth century. The postwar Japanese case instead captures a formative period in which the welfare state itself became a direct lender to low-income households. The results suggest that the state operated not only as a provider of transfers and insurance but also as a creditor that distinguished between households deemed capable of repayment and those left to rely on pawnshops, informal networks, or consumption compression.

The transformation documented here was not a simple transition from pawnshops to modern public finance. Welfare lending did not eliminate collateral-based credit; instead, it added a new lending technology based on assessments of household rehabilitation capacity. Postwar financial modernization involved not institutional unification but the hierarchical allocation of low-income households across institutions employing different standards of creditworthiness.

\newpage

\section{Institutional Background: Economic Growth and the Reorganization of Credit for Low-Income Households}

\subsection{Economic growth and persistent vulnerability}

From the late 1950s through the early 1960s, the Japanese economy expanded at an extraordinary pace. Using the Cabinet Office's historical national accounts series, real GDP, indexed to 100 in 1955, rose to approximately 152 in 1960, 170 in 1961, and 237 in 1965 \citep{CabinetOffice2000}. In other words, the real size of the economy increased by roughly 2.4 times within a decade.

Reliance on public assistance also declined during this period. The number of public-assistance recipients per 1,000 population fell from 21.6 in fiscal year 1955 to 16.3 in fiscal year 1965. Yet the decline in public assistance did not mean that low-income insecurity had disappeared. The 1958 \emph{White Paper on Health and Welfare} estimated that 2.4 million of Japan's 18.91 million households---12.7 per cent---belonged to the category of ``low-consumption-level households.'' The same report stressed that the gains from rising average consumption were distributed unevenly across social strata \citep{MinistryofHealthandWelfare1958}.

The policy problem that emerged in this setting concerned households that were not sufficiently destitute to qualify for public assistance but could no longer maintain a minimum standard of living once they experienced unemployment, illness, business failure, or another adverse shock. Contemporary welfare administration referred to such households as the \emph{borderline stratum} (\emph{b\=odārain-s\=o}): households situated between public-assistance recipients and economically secure households \citep{Sakai2023}.

This position was institutionally important. Public assistance provided transfers to households already unable to maintain minimum living standards from their income and assets. Policies for borderline households, by contrast, sought to prevent descent into public assistance by temporarily financing business investment, medical treatment, job search and preparation, housing, or education. During the high-growth era, low-income policy thus expanded beyond ex post relief toward a preventive policy based partly on the provision of credit.

\subsection{The emergence of welfare credit}

The institutional foundation of this preventive approach was the Household Rehabilitation Movement (\emph{setai k\=osei und\=o}) led by local welfare commissioners (\emph{minsei-iin}). The 1952 National Conference of Welfare and Child Welfare Commissioners adopted a resolution promoting the movement, under which welfare commissioners provided continuing advice and assistance to low-income households. In fiscal year 1955, the central government began subsidising the Household Rehabilitation Fund Loan Program (\emph{setai k\=osei shikin kashitsuke seido}), which became the movement's principal financial instrument \citep{MinistryofHealthandWelfare1958}. A separate medical-expense loan program was introduced in fiscal year 1957.

Loans were administered by prefectural social welfare councils. Welfare commissioners identified and advised households, mediated applications, and provided guidance concerning both household management and repayment. A major revision in April 1961 expanded the available loan categories, increased lending ceilings, lengthened repayment periods, and incorporated medical-expense loans into the Household Rehabilitation Fund as medical-treatment loans. By the end of fiscal year 1960, cumulative lending had reached 3.934 billion yen to 109,851 borrowers, while approval rates generally remained around 80 per cent of applications \citep{MinistryofHealthandWelfare1961}.

The program did not distribute unrestricted cash uniformly among low-income households. Rather, it provided purpose-specific loans for livelihood businesses, job preparation, vocational training, housing, education, medical treatment, and related forms of household rehabilitation. Its objective was not simply to fill a current income deficit but to finance expenditures expected to restore or strengthen a household's future earning capacity. The Ministry of Health and Welfare described the combination of finance and individual guidance by welfare commissioners as a defining feature of the program \citep{MinistryofHealthandWelfare1964}.

This design contained a fundamental tension. Eligible households had to be poor enough that they could not obtain the required funds from ordinary financial institutions. At the same time, because assistance took the form of a loan rather than a transfer, they also had to appear capable of future repayment.

Contemporary legislators recognised this contradiction clearly. During a 1964 House of Representatives debate concerning welfare loans for fatherless households, Hasegawa Tamotsu argued that ``the people who need [the loans] most are precisely those who look too risky to repay.'' His concern was that the households with the greatest need might have the weakest apparent repayment capacity and therefore be screened out of the program. Government representatives acknowledged that a high repayment rate was not, by itself, sufficient evidence of successful administration. The debate also raised the possibility that welfare commissioners or local offices discouraged risky households from applying before a formal rejection could be recorded. The Ministry replied that it had not investigated such pre-application screening \citep{NationalDietLibrary1964}.

Although the immediate subject of this debate was the loan program for fatherless households, the underlying tension applied more broadly to welfare credit: the state sought to lend to poor households while preserving the revolving public fund through repayment. Welfare-loan users may therefore have differed from the poorest households. In addition to a specific financial need, they may have possessed labour capacity, a viable business or rehabilitation plan, family resources, or a minimum material basis from which recovery appeared possible.

\subsection{Welfare loans and public pawnshops}

Before the expansion of welfare lending, pawnshops had long supplied small loans to low-income households excluded from bank credit. Public pawnshops (\emph{k\=oeki shichiya}) were operated by municipalities or social welfare corporations and offered livelihood and business funds on terms more favourable than those of private pawnshops. In 1957, the standard monthly interest rate at public pawnshops was 3 per cent, compared with approximately 9 per cent at private pawnshops. Public pawnshops also calculated interest in half-month units, allowed longer periods before forfeiture, and returned any surplus remaining after the sale of forfeited collateral. Nevertheless, there were only 779 public pawnshops nationwide, compared with 21,224 private establishments \citep{MinistryofHealthandWelfare1958}.

By March 1961, 848 public pawnshops were operating nationwide. During fiscal year 1960, they made approximately 2.24 million loans with a total value of 3.14 billion yen. The average loan was only about 1,400 yen. Wage earners accounted for approximately 60 per cent of users and the self-employed for another 20 per cent. The combination of very small average loan sizes and extremely large transaction volumes suggests that public pawnshops continued to provide short-term liquidity for ordinary household needs \citep{MinistryofHealthandWelfare1961}.

Placing the national trends in welfare lending and public pawnbroking side by side illustrates a reorganisation of low-income finance between the late 1950s and early 1960s.

\begin{table}[h]
\centering
\caption{Welfare lending and public pawnshops, 1959--1963}
\label{tab:welfareandpawn}
\begin{tabular}{lccc}
\toprule
 & 1959 & 1961 & 1963 \\
\midrule
\multicolumn{4}{l}{\textbf{Household Rehabilitation Fund}} \\
Households receiving new loans & 27,386 & 30,673 & 31,812 \\
Total lending (billion yen, nominal) & 0.973 & 1.590 & 2.157 \\
Average amount per household (yen) & 35,511 & 51,828 & 67,815 \\
\midrule
\multicolumn{4}{l}{\textbf{Public pawnshops}} \\
Number of establishments & 836 & 831 & 756 \\
Number of loans (millions) & 2.499 & 1.903 & 1.476 \\
Total lending (billion yen, nominal) & 3.294 & 3.095 & 3.005 \\
Average amount per loan (yen) & 1,318 & 1,626 & 2,036 \\
\bottomrule
\end{tabular}
\end{table}

\noindent\emph{Source:} Ministry of Health and Welfare (1964), Tables 2-10-6 and 2-10-8.

Between 1959 and 1963, the annual amount lent through the Household Rehabilitation Fund increased by approximately 122 per cent, while the average amount per borrowing household rose from roughly 35,500 yen to 67,800 yen. Over the same period, the number of public-pawnshop loans declined by approximately 41 per cent. Total pawnshop lending, however, remained close to 3 billion yen, and the average amount per loan rose by approximately 54 per cent.

These trends do not establish that welfare lending simply displaced public pawnshops. They instead suggest a restructuring in which welfare lending expanded as a source of relatively large, purpose-specific rehabilitation finance, while public pawnshops continued to conduct a very large number of small transactions. The Ministry itself attributed the decline in pawnshop activity partly to the expansion of social protection and instalment sales, but it also emphasised that public pawnshops remained relevant to the everyday lives of low-income households \citep{MinistryofHealthandWelfare1963, MinistryofHealthandWelfare1964}.

An important reason for this coexistence was that the Household Rehabilitation Fund did not comprehensively finance everyday consumption. The 1964 \emph{White Paper on Health and Welfare} explained the stagnation of subsistence loans partly by noting that loans for living expenses could not be made independently of another qualifying purpose. At the same time, lending for business rehabilitation and education---expenditures expected to generate future income or benefits---was growing rapidly \citep{MinistryofHealthandWelfare1964}.

Obtaining a welfare loan therefore did not necessarily satisfy all of a household's financial needs. Business loans were earmarked for starting or rebuilding an enterprise, medical-treatment loans for treatment expenses, and educational loans for schooling. Food, rent, utilities, and family consumption during illness or business reconstruction often had to be financed separately.

Evidence from actual borrowers supports this interpretation. A reconstructed survey of 654 households that received medical-expense loans in Kanagawa between 1956 and 1960 shows that 449 cases, or 68.7 per cent, reported that the loan was insufficient to cover treatment costs. In 233 cases---36 per cent of the total---the remaining cost was financed through additional borrowing \citep{Sakai2023}.

The relationship between welfare loans and other credit sources was therefore not necessarily one of simple substitution. A low-income household might finance business investment, medical treatment, or education through a welfare loan; borrow from an employer, relatives, friends, or neighbours to cover everyday consumption; postpone payments through purchases on account; and pawn household goods to meet immediate cash shortages. Rather than graduating from older or informal credit after obtaining public finance, households may have combined multiple sources with different purposes and maturities---a pattern of \emph{layered borrowing} or \emph{credit stacking}.

Pawnshop use itself was not necessarily a one-off emergency response. Using loan data from more than 250 municipalities, researchers find that the increase in pawnshop lending during the 1918--20 influenza pandemic was driven principally by larger borrowing among regular pawnshop customers rather than by a substantial influx of new borrowers \citep{Inoue2021}. Pawnshops could therefore serve as repeatedly used liquidity providers for households facing persistent financial constraints.

For these reasons, the decline in the number of public-pawnshop transactions cannot be interpreted as evidence that welfare loans directly displaced pawnshops. Economic growth, the extension of social insurance, the diffusion of durable consumer goods, and the development of instalment credit were occurring simultaneously. Moreover, welfare loans and pawnshops could be used by the same households at different times and for different purposes.

\subsection{Two technologies of public credit and layered borrowing}

The Household Rehabilitation Fund and public pawnshops both supplied credit to low-income households that could not obtain sufficient funds in ordinary financial markets, but they did so through different lending technologies. The Household Rehabilitation Fund pursued poverty prevention and household rehabilitation through purpose-specific loans for business, medical treatment, housing, and education. Prefectural social welfare councils made the loans, while welfare commissioners mediated access, assisted households, and provided post-loan guidance. The Ministry of Health and Welfare explicitly treated the combination of finance and individual supervision as a central feature of the program \citep{MinistryofHealthandWelfare1964}.

Welfare lending therefore depended not only on current income but also on the purpose of the loan and the prospect that the financed expenditure would contribute to household recovery. A livelihood-business loan presupposed some degree of business feasibility; a medical-treatment loan was connected to recovery and a return to ordinary life or employment; and an educational loan was intended to sustain schooling and improve future employment prospects. Welfare credit can thus be understood as credit based partly on the household's circumstances, financial plan, and expected capacity for rehabilitation. The available data, however, do not reveal the actual deliberations behind individual approval or rejection decisions.

Public pawnshops relied on a different basis of credit. The Ministry described them as institutions that supplied low-income households with ``simple and rapid'' access to livelihood or business finance. Although public pawnshops operated under administrative rules and identity requirements, the value of the pledged object ultimately secured the loan. Welfare lending evaluated the possible rehabilitation of a household; public pawnbroking converted currently held physical assets into liquidity \citep{MinistryofHealthandWelfare1958, MinistryofHealthandWelfare1963}.

The difference between these technologies did not make them mutually exclusive. Welfare loans were purpose-specific and did not comprehensively insure everyday living expenses. Medical-treatment loans could help pay hospital bills without covering rent, food, utilities, or family consumption during hospitalisation. Business loans could finance premises or equipment without sustaining household consumption until the enterprise began to generate revenue.

The Kanagawa borrower evidence makes this limitation concrete: 68.7 per cent of medical-loan cases reported that the loan did not cover treatment costs, and 36 per cent financed the shortfall through further debt \citep{Sakai2023}. Welfare lending may therefore have coexisted with pawning, employer loans, borrowing from relatives or neighbours, and purchases on account. Low-income households could combine public loans for designated expenditures with other forms of credit for ordinary consumption and immediate liquidity. This is more appropriately understood as \emph{layered borrowing} than as graduation from older credit institutions.

Pawnshops, similarly, were not necessarily used only once in an exceptional emergency. Evidence from the influenza pandemic shows that established users intensified their borrowing during a health shock \citep{Inoue2021}. If past public-pawnshop users in postwar Japan also continued to pawn goods, this would suggest that pawnbroking retained a role as a repeatedly used source of small-scale liquidity even after welfare lending had emerged. If welfare-loan users continued to borrow from employers, friends, or neighbours, public credit should be understood as an addition to, rather than a full replacement for, relational and informal finance.

This institutional background generates two distinct questions for the empirical analysis. The first concerns \textbf{selection across institutions}. What kinds of households obtained welfare loans, and what kinds used public pawnshops? If welfare lending placed weight on rehabilitation and repayment capacity, its users need not have been the most destitute households. They may instead have possessed a minimum employment or asset base and a concrete prospect of household recovery. Public pawnshops, by contrast, may have served households facing immediate liquidity shortages associated with unemployment, illness, or other shocks and able to obtain credit primarily by pledging goods.

The second concerns \textbf{substitution versus layered use}. Did welfare-loan users cease relying on pawnshops and neighbourhood borrowing, or did they continue to combine multiple financial instruments? Continued pawning among past public-pawnshop users would suggest that public pawnbroking was not merely a one-time emergency device but a recurrent source of liquidity. Continued borrowing from employers, friends, or neighbours among welfare-loan users would indicate that public lending did not fully displace informal finance but was layered on top of existing sources of credit.

\newpage

\section{Data and Sample}

This paper draws on the 1961 \emph{Kanagawa Survey of the Borderline Poor} (\emph{Kanagawa-ken Kyōdō Chōsa Shiryō}), a restored dataset of 6,152 low-income households interviewed in six Kanagawa prefecture cities: Yokohama, Kawasaki, Yokosuka, Kamakura, Hiratsuka, and Chigasaki. The survey was conducted by the Kanagawa Social Welfare Council between December 1960 and March 1961 as part of a systematic assessment of households in the borderline poverty stratum. The dataset records detailed information on household composition, employment, income sources, asset holdings, causes of low income, use of public assistance, welfare loans, and public pawnshops.

The analysis covers the full population of surveyed households: 6,152 families representing the sample frame of the survey at the time of enumeration. The survey records the principal breadwinner's occupation, age, health status, and disability status; household size and composition; monthly and annual income; identified causes of low income (unemployment, business failure, declining earnings, illness, disability, war damage); and institutional use histories. For each household, the survey records whether it had ever received public assistance, used any welfare loans (maternal welfare, household rehabilitation, medical, or educational), or pawned goods at public pawnshops. The survey also documents household coping methods when living expenses were insufficient: borrowing from friends or neighbors, employer borrowing, purchases on account, reducing food consumption, withdrawing savings, and selling assets.

The key outcome variables are binary indicators of welfare-loan use (any welfare loan ever received) and public-pawnshop use (any history of pawning goods). The sample includes households in both categories as well as those using neither institution, enabling comparison of institutional allocation patterns within the low-income population. The independent variables of central interest are risk indicators reflecting causes of low income (business failure, prolonged illness, disability, unemployment, war damage, income decline) and household economic position (public assistance history, reported livelihood capacity, principal breadwinner health status, number of household assets). Income is recorded both at the household level and disaggregated by source.

Descriptive statistics appear in Table 1. Among the 6,152 households, 307 (5.0\%) had received any welfare loan, 110 (1.8\%) disability medical-care support, 115 (1.9\%) maternal welfare funds, 176 (2.9\%) educational loans, and 187 (3.0\%) medical-expense loans. Public pawnshop use was more common: 378 households (6.1\%) reported ever pawning goods. Only 30 households (0.5\%) had used both welfare loans and public pawnshops, indicating that the two institutions served largely distinct populations. A history of public assistance receipt was present in 1,539 households (25.0\%). The average household reported 3.2 assets (range 0--12), and the principal breadwinner was recorded as in poor health in 856 households (13.9\%).

\begin{table}[!h]
\centering
\small
\caption{Summary Statistics by Institutional-Use Group}
\label{tab:summary}
\begin{tabular}{lrrrr}
\toprule
Characteristic & Neither & Welfare & Pawnshop & Both \\
\midrule
\multicolumn{5}{l}{\textbf{Demographics}} \\
N & 5,045 & 729 & 348 & 30 \\
Head age (years) & 50.6 & 48.5 & 50.2 & 49.8 \\
Household size & 3.9 & 4.1 & 3.8 & 3.9 \\
Workers & 1.3 & 1.4 & 1.2 & 1.3 \\
Female head (\%) & 6.3 & 7.5 & 6.6 & 6.7 \\
\midrule
\multicolumn{5}{l}{\textbf{Economic Indicators}} \\
Asset count & 2.8 & 3.5 & 2.6 & 2.9 \\
Public assistance (\%) & 20.4 & 17.7 & 15.8 & 26.7 \\
Low livelihood (\%) & 29.1 & 26.0 & 35.1 & 36.7 \\
Poor health (\%) & 12.8 & 12.1 & 15.5 & 13.3 \\
\midrule
\multicolumn{5}{l}{\textbf{Risk Indicators}} \\
Business failure (\%) & 7.7 & 15.3 & 6.6 & 10.0 \\
Prolonged illness (\%) & 16.8 & 22.1 & 18.4 & 16.7 \\
Unemployment (\%) & 10.3 & 10.9 & 14.7 & 13.3 \\
Income decline (\%) & 13.1 & 13.6 & 19.3 & 13.3 \\
War damage (\%) & 8.9 & 8.8 & 7.2 & 3.3 \\
\bottomrule
\end{tabular}

\noindent\small\emph{Note:} N=6,152 households. Continuous vars: means. Binary vars: percentages. Groups: Neither (no institutional credit), Welfare (welfare only), Pawnshop (pawnshop only), Both (both institutions).
\end{table}

\section{Empirical Strategy}

The analysis estimates the probability of welfare-loan use and public-pawnshop use among low-income households using linear probability models with heteroskedasticity-consistent standard errors (HC3). The general specification is:

\begin{equation}
Y_i = \alpha + \beta X_i + \varepsilon_i,
\end{equation}

\noindent where $Y_i$ is a binary indicator of institutional use (welfare loan or public pawnshop), $X_i$ is a vector of household characteristics including risk indicators and economic position measures, and $\varepsilon_i$ is the error term. Standard errors are computed using the HC3 estimator, which is robust to arbitrary heteroskedasticity and performs better in finite samples than HC0 or HC1 (Davidson and MacKinnon 2004).

The linear probability model (LPM) is chosen for two reasons. First, it facilitates interpretation of coefficients as marginal effects on the probability of institutional use, measured in percentage points. Second, LPM estimation is more robust to model misspecification when the conditional probability is bounded away from zero and one; in our sample, welfare-loan prevalence is 5\% and pawnshop prevalence is 6\%, providing sufficient variation to make LPM estimates interpretable. Logit models are estimated as robustness checks and reported alongside LPM coefficients in the results tables.

Model specifications proceed from parsimonious to more complete. The main specification controls for household income (log annual income), geographic fixed effects (five dummy variables for cities), household size (number of members), principal breadwinner employment status (self-employed, wage-employed, unemployed, other), and risk indicators including binary variables for business failure, prolonged illness, disability, unemployment, income decline, and war damage. A second specification adds measures of household economic position: public assistance history (binary indicator), reported livelihood capacity (binary indicator of low capacity), principal breadwinner health status (binary indicator of poor health), and asset count (number of household assets).

This empirical strategy is descriptive, not causal. The analysis documents correlations between household characteristics and institutional use, not the treatment effects of receiving credit. The survey records historical use of institutions, not dates of application, approval, or rejection; accordingly, the results describe sorting patterns rather than the causal effects of borrowing. Reverse causality (borrowing could affect observed outcomes) is therefore not a concern; instead, the results reflect how household characteristics at a point in time are associated with patterns of institutional use.

\section{Main Results}

\subsection{Welfare lending and the nature of shocks}

Table 2 reports the association between risk factors and household characteristics with the probability of welfare-loan use. Even after controlling for income, location, household composition, and employment, specific risk factors are strongly associated with welfare-loan use. Business failure increases the probability of welfare-loan use by 8.6 percentage points (SE = 1.2, LPM coefficient = 0.0856), a substantial effect given a baseline prevalence of 5.0\%. Prolonged illness is associated with a 5.1-percentage-point increase (SE = 1.1, LPM coefficient = 0.0506), while disabled household members increase the probability by 4.6 percentage points (SE = 1.9). These associations are robust to model specification and significant at conventional levels.

Conversely, low livelihood capacity is associated with a 3.2-percentage-point lower probability of welfare-loan use (SE = 1.9). A history of public assistance is associated with a 6.1-percentage-point reduction in welfare-loan probability (SE = 1.2). Each additional household asset is associated with a 0.54-percentage-point increase (SE = 0.1). This pattern is consistent with welfare lending targeting specific, addressable problems amenable to targeted rehabilitation finance (business investment, medical treatment) among households with economic capacity to repay, rather than operating as unrestricted support for the destitute.

The individual welfare loan programmes correspond closely to their stated purposes. Use of the Household Rehabilitation Fund is associated with self-employment and demand for business capital. Educational lending is significantly more prevalent among households with school-age members. Medical lending is concentrated among households reporting prolonged illness or reporting medical expenses as a cause of low income. This purpose-specific targeting indicates that welfare loans functioned as earmarked credit tied to particular rehabilitation plans, not as general-purpose poverty relief.

\begin{table}[!h]
\centering
\small
\caption{Risk Factors and Welfare-Loan Use}
\label{tab:shocks}
\begin{tabular}{lrr}
\toprule
Variable & LPM & Logit AME \\
\midrule
Business failure & $0.0856^{***}$ & $0.0830^{***}$ \\
 & (0.0120) & (0.0120) \\
Prolonged illness & $0.0506^{***}$ & $0.0494^{***}$ \\
 & (0.0110) & (0.0110) \\
Disabled household member & $0.0464^{**}$ & $0.0434^{**}$ \\
 & (0.0190) & (0.0190) \\
Low living ability & $-0.0323$ & $-0.0353$ \\
 & (0.0190) & (0.0190) \\
Public assistance & $-0.0613^{***}$ & $-0.0592^{***}$ \\
 & (0.0120) & (0.0120) \\
Asset count & $0.0054^{***}$ & $0.0057^{***}$ \\
 & (0.0010) & (0.0010) \\
\midrule
Observations & 6,152 & 6,152 \\
\midrule
\multicolumn{3}{p{0.4\textwidth}}{\emph{Controls:} Log income, city FE, household size, breadwinner employment.} \\
\bottomrule
\end{tabular}

\noindent\small\emph{Note:} Dependent variable: welfare-loan use. HC3-robust SE in parentheses. $^{*}p<0.10$, $^{**}p<0.05$, $^{***}p<0.01$. Sample: 6,152 households.
\end{table}

Program-specific targeting is further evident in Table 5, which reports the distribution of users and targeting accuracy across welfare programs. The Household Rehabilitation Fund reached 307 households (4.99\% of the sample) with a targeting accuracy of 13.4\%, meaning that 13.4\% of the program's users had experienced business failure, its primary target. The maternal welfare fund showed the highest targeting accuracy at 92.2\%, concentrated among 115 beneficiaries (1.87\%). Educational loans reached 176 households (2.86\%) with targeting accuracy of 31.0\%. Medical-expense loans, serving 187 households (3.04\%), had targeting accuracy of 16.6\%. These figures demonstrate that welfare programs were deliberately aimed at specific needs, even if each program captured a relatively modest share of the low-income population.

\begin{table}[!h]
\centering
\small
\caption{Targeting Accuracy by Program}
\label{tab:targeting}
\begin{tabular}{lrrrr}
\toprule
Program & N Users & \% Sample & Accuracy & Risk Ratio \\
\midrule
Public assistance & 1,539 & 25.0 & 62.5 & 500.41 \\
Disability medical & 110 & 1.8 & 30.0 & 4.02 \\
Household Rehab Fund & 307 & 5.0 & 13.4 & 1.80 \\
Maternal welfare & 115 & 1.9 & 92.2 & 3.15 \\
Medical expense loan & 187 & 3.0 & 16.6 & 1.79 \\
Fatherless children & 28 & 0.5 & 44.6 & 1.49 \\
Special education & 176 & 2.9 & 31.0 & 1.05 \\
Public pawnshop & 378 & 6.1 & 9.1 & 1.24 \\
\bottomrule
\end{tabular}

\noindent\small\emph{Note:} Accuracy = \% of users with primary risk indicator. Risk Ratio = ratio of risk among users to general population. Columns 3-4 in percentages.
\end{table}

\subsection{Economic position and institutional allocation}

A striking pattern emerges when comparing welfare-loan users to public-pawnshop users: the clearest distinction lies not in the type of shock experienced, but in the household's underlying economic and material position. Table 3 presents results from models restricting the sample to households using either welfare loans or public pawnshops, but not both. Among this subset, a history of public assistance is associated with a 15.6-percentage-point lower probability of belonging to the welfare-loan group (SE = 3.2), a large and highly significant effect. Households recorded by the survey as having low livelihood capacity are 14.5 percentage points less likely to use welfare loans (SE = 3.8), and households whose principal breadwinner was in poor health are 11.4 percentage points less likely (SE = 3.5).

By contrast, each additional household asset is associated with a 3.8-percentage-point higher probability of welfare-loan use (SE = 0.8). This pattern holds across multiple model specifications and is robust to including controls for income level. The effect is economically significant: a household with five assets (compared to one) would have a 15-percentage-point higher probability of welfare-loan use, a substantial difference.

These results are consistent with welfare lending depending on the household's prospective capacity for rehabilitation and repayment, whereas public pawnshops relied on currently available physical collateral. Households with a public-assistance history were likely perceived by welfare officials as unlikely to repay or to successfully rehabilitate even with targeted support. Households with low livelihood capacity and poor health status faced similar barriers to welfare-loan approval. Households with more assets, by contrast, signaled an economic foundation from which recovery might be possible. The pattern does not reflect simple economic need (income controls are included in all models) but rather the credit-assessment criteria implicit in welfare lending.

\begin{table}[!h]
\centering
\small
\caption{Economic Position and Institutional Choice}
\label{tab:position}
\begin{tabular}{lr}
\toprule
Variable & Coefficient \\
\midrule
Public assistance history & $-0.1560^{***}$ \\
 & (0.0320) \\
Low livelihood capacity & $-0.1450^{***}$ \\
 & (0.0380) \\
Breadwinner poor health & $-0.1140^{***}$ \\
 & (0.0350) \\
Asset count (per item) & $0.0380^{***}$ \\
 & (0.0080) \\
\midrule
Observations & 725 \\
\midrule
\multicolumn{2}{p{0.35\textwidth}}{\emph{Controls:} Log income, city FE, household size, breadwinner employment.} \\
\bottomrule
\end{tabular}

\noindent\small\emph{Note:} Dep. var.: welfare-loan use (vs. pawnshop-only). Sample: 725 households. HC3-robust SE in parentheses. $^{*}p<0.10$, $^{**}p<0.05$, $^{***}p<0.01$.
\end{table}

\subsection{Credit layering and institutional use}

A third finding is that welfare lending did not fully replace older forms of borrowing. Table 4 documents the coping strategies and credit sources used by welfare-loan users and public-pawnshop users compared to non-users of institutional credit. Welfare-loan users were 10.3 percentage points more likely to report pawning goods when living expenses were insufficient (SE = 2.1), compared to non-users. Households with a history of public-pawnshop use were 38 percentage points more likely to report current pawning (SE = 3.2), a very large effect indicating that past pawnshop use predicts continued reliance on this institution.

Borrowing from friends and neighbors was more common among welfare-loan users (8.5 pp higher probability; SE = 2.3) but also more common among households with public-pawnshop histories (7.2 pp; SE = 2.9). Employer borrowing was slightly more common among pawnshop users (3.1 pp; SE = 1.8). Consumption compression strategies—reducing food intake, withdrawing savings, and selling assets—were used across all groups but were actually slightly less common among welfare-loan users, suggesting that access to institutional credit did provide some degree of consumption smoothing.

This evidence indicates that households did not graduate from older credit sources after obtaining welfare loans, but rather layered multiple instruments. The result is consistent with welfare loans serving purpose-specific needs (medical care, business investment, education) while other borrowing provided general liquidity for living expenses. Public pawnshops, similarly, were not one-time emergency responses but repeatedly used sources of small-scale credit for the persistently financially constrained.

\begin{table}[!h]
\centering
\small
\caption{Coping Strategies by Institutional Access}
\label{tab:coping}
\begin{tabular}{lrrrr}
\toprule
Coping Method & None & Welfare & Pawnshop & Both \\
\midrule
Pawning goods & 5.2 & 10.3 & 44.3 & 50.0 \\
Employer borrowing & 14.5 & 16.3 & 27.0 & 20.0 \\
Friend/neighbor borrow & 18.8 & 23.6 & 29.0 & 50.0 \\
Asset sale & 2.7 & 1.5 & 4.6 & 0.0 \\
Savings withdrawal & 5.5 & 2.9 & 2.3 & 0.0 \\
Food compression & 25.8 & 22.9 & 27.6 & 20.0 \\
\midrule
N & 5,045 & 729 & 348 & 30 \\
\bottomrule
\end{tabular}

\noindent\small\emph{Note:} Percentages within institutional-access groups. None (N=5,045), Welfare (N=729), Pawnshop (N=348), Both (N=30). Strategies reported when insufficient living expenses. No controls; observed shares.
\end{table}

The intensity of layered borrowing is further documented in Table 8, which reports the distribution of households by number of institutional credit sources used. The majority of households (5,045, or 82.0\%) used no institutional credit from either welfare loans or public pawnshops. Among those using institutional credit, 1,077 households (17.5\%) relied on a single source, while only 30 households (0.5\%) used both welfare loans and public pawnshops simultaneously. This small overlap confirms that the two institutions largely served distinct populations rather than overlapping user groups. Among households using institutional credit, the mean number of coping strategies employed was 0.96 for single-source users and 1.40 for those combining both institutions, suggesting that credit diversification is indeed associated with additional reliance on non-institutional coping methods.

\begin{table}[!h]
\centering
\small
\caption{Credit Intensity and Coping Strategy Use}
\label{tab:intensity}
\begin{tabular}{lrrrr}
\toprule
Institutional Sources & N & \% & Coping Strats & Food Comp \\
\midrule
No institutional & 5,045 & 82.0 & 0.72 & 25.8 \\
One source & 1,077 & 17.5 & 0.96 & 24.4 \\
Both sources & 30 & 0.5 & 1.40 & 20.0 \\
\bottomrule
\end{tabular}

\noindent\small\emph{Note:} Coping Strats = mean number of strategies when insufficient expenses. Food Comp = \% reducing food. \% in columns 3-4 are percentages.
\end{table}

\section{Robustness Checks}

The main results are robust to alternative specifications and sampling procedures. Leave-one-city-out analysis re-estimates the main model excluding each city in turn. Table 6 presents the leave-one-city-out results across all six cities. The coefficient on business failure ranges from 0.0725 to 0.1046 across city omissions, with a mean of 0.0825, nearly identical to the full-sample estimate of 0.0856. The coefficient on prolonged illness ranges from 0.0363 to 0.0887 with a mean of 0.0534, very close to the full-sample estimate of 0.0506. Coefficients on public-assistance history range from $-0.0453$ to $-0.0728$ with a mean of $-0.0613$. The coefficient on low living ability shows variation, ranging from $-0.0248$ to $-0.0380$ with a mean of $-0.0325$. This stability across city subsamples suggests that institutional sorting patterns were consistent across the Kanagawa cities surveyed.

\begin{table}[!h]
\centering
\small
\caption{Leave-One-City-Out Robustness}
\label{tab:loco}
\begin{tabular}{lrrrrrr}
\toprule
 & C1 & C2 & C3 & C41 & C42 & C43 \\
\midrule
\multicolumn{7}{l}{\textbf{Business Failure}} \\
Coeff. & 0.0825 & 0.1046 & 0.0840 & 0.0795 & 0.0857 & 0.0726 \\
SE & (0.0275) & (0.0228) & (0.0226) & (0.0197) & (0.0198) & (0.0196) \\
\midrule
\multicolumn{7}{l}{\textbf{Prolonged Illness}} \\
Coeff. & 0.0887 & 0.0363 & 0.0482 & 0.0512 & 0.0467 & 0.0513 \\
SE & (0.0266) & (0.0176) & (0.0183) & (0.0167) & (0.0165) & (0.0167) \\
\midrule
\multicolumn{7}{l}{\textbf{Public Assistance}} \\
Coeff. & $-0.0453$ & $-0.0570$ & $-0.0728$ & $-0.0625$ & $-0.0634$ & $-0.0586$ \\
SE & (0.0154) & (0.0104) & (0.0103) & (0.0095) & (0.0094) & (0.0095) \\
\midrule
\multicolumn{7}{l}{\textbf{Low Living Ability}} \\
Coeff. & $-0.0248$ & $-0.0310$ & $-0.0380$ & $-0.0346$ & $-0.0297$ & $-0.0339$ \\
SE & (0.0151) & (0.0105) & (0.0109) & (0.0097) & (0.0098) & (0.0097) \\
\bottomrule
\end{tabular}

\noindent\small\emph{Note:} Each model excludes one city. All models control: log income, household size, breadwinner employment, city FE. HC3-robust SE in parentheses. Sample sizes: C1 N=2919, C2 N=4802, C3 N=4986, C41 N=6008, C42 N=5959, C43 N=5981.
\end{table}

Table 7 reports results across alternative model specifications, comparing the fit and structure of different estimation approaches. The minimal specification (controlling only for city fixed effects) yields an R-squared of 0.0034. Adding demographic controls (household head age, household size, number of workers) increases R-squared to 0.0168. The full specification, including all risk indicators and household economic-position measures, yields an R-squared of 0.0384. A specification using a composite risk index rather than individual risk indicators yields an R-squared of 0.0203, lower than the full specification, suggesting that individual risk factors carry distinct information relevant to institutional allocation. These results confirm that while the explained variance remains modest (as is typical in binary choice models with many households), the specific risk indicators and economic-position variables consistently improve model fit.

\begin{table}[!h]
\centering
\small
\caption{Model Specification Comparison}
\label{tab:specs}
\begin{tabular}{lrrr}
\toprule
Specification & Model 1 & Model 2 & Model 3 \\
\midrule
R-squared & 0.0034 & 0.0168 & 0.0384 \\
Observations & 6,131 & 6,131 & 6,131 \\
\midrule
\multicolumn{4}{l}{\emph{Controls:}} \\
City FE & Yes & Yes & Yes \\
Income & No & No & Yes \\
Household size & No & Yes & Yes \\
Breadwinner employment & No & Yes & Yes \\
Risk indicators & No & No & Yes \\
Economic position & No & No & Yes \\
\bottomrule
\end{tabular}

\noindent\small\emph{Note:} Model 1: city FE only. Model 2: demographics (age, size, employment). Model 3: full specification with all risk and position controls. HC3-robust estimates.
\end{table}

Logit models with average marginal effects yield results similar to LPM coefficients. For business failure, logit AME is 0.0830 compared to LPM 0.0856. For prolonged illness, logit AME is 0.0494 compared to LPM 0.0506. These similarities confirm that the linear model provides a reasonable approximation in this range.

Alternative specifications restricting to households experiencing business failure show that public-assistance history remains strongly negative even among households facing the same shock. This reinforces the interpretation that welfare-lending decisions depended on household rehabilitation capacity, not merely need.

\section{Institutional Allocation and the Welfare State}

The results characterize how low-income households were sorted across credit institutions when the welfare state entered the credit market. Welfare lending was not extended uniformly to the poorest households, but concentrated among those with specific financing needs and economic foundations suggesting repayment capacity. Public pawnshops served as rapid-access credit based on physical collateral. This institutional differentiation was not accidental but reflected deliberate design: welfare loans required applications and bureaucratic review; pawnshops required only an asset.

The result was hierarchical allocation. At the top sat households with identifiable needs and minimum economic foundations---they accessed welfare loans. Below were households with immediate liquidity needs and collateral---they relied on pawnshops. At the bottom were the poorest, those with public-assistance histories and minimal assets---largely excluded from institutional credit. This stratification reflects the inherent tension in welfare lending: providing credit while maintaining creditworthiness standards.

\section{Conclusion}

This paper examined how low-income households in postwar Japan were allocated across welfare lending and public pawnshops. Three main findings emerge. First, welfare lending targeted specific problems: business failure increased welfare-loan probability by 8.6 percentage points, prolonged illness by 5.1 percentage points. Second, the clearest distinction between users lay in economic position: public-assistance histories decreased probability by 15.6 percentage points, while household assets increased it by 3.8 percentage points per item. Third, welfare lending did not displace older borrowing; households layered institutional credit with informal borrowing and consumption compression.

These patterns reflect institutional design. Welfare loans required applications and rehabilitation assessment. Pawnshops required collateral. The welfare state therefore created new boundaries of creditworthiness within the poor population. Credit access was determined not by need alone, but by perceived capacity for repayment.

Future research should examine whether welfare-loan recipients experienced improved long-term outcomes. Understanding how multiple institutions coexist and allocate access in constrained credit markets remains important for development policy. The postwar Japanese case demonstrates that state entry into credit markets, even with poverty-reduction goals, shapes access through administrative mechanisms and market procedures.

\newpage

\bibliographystyle{abbrvnat}
\bibliography{references}

\end{document}
